\section{Reproducibility}\label{app:repro}
All runs write a fully resolved configuration, the git commit, an environment snapshot, a
per-round metrics log and a 1\,Hz power and memory log. Models and datasets are pinned to Hub
commits; the data manifest records file and id hashes of every split and evaluation set. Each
table in this paper is regenerated from those files by \texttt{scripts/make\_paper\_tables.py}.
The experiment grids are YAML files expanded into one resolved configuration per run and executed
by a sequential, resumable queue (one run at a time; finished runs are skipped and an interrupted
run resumes from its own state file), so the full set of runs can be re-executed with one command.

\section{Scoring backends and throughput}\label{app:backends}
% generated by scripts/make_paper_tables.py from runs/agreement/S2_vllm_vs_hf_s0/summary.json; do not edit
\begin{table}[t]
  \centering\small
  \caption{Teacher-scoring backends agree: vLLM \texttt{prompt\_logprobs} vs.\ a plain HF forward pass (S2 as stand-in teacher, 20 sequences).}\label{tab:agreement}
  \begin{tabular}{lr}
    \toprule
    Metric & Value \\
    \midrule
    Tokens compared & 10{,}240 \\
    Mean $|\Delta|$ (nats) & 0.0152 \\
    Median $|\Delta|$ & 0.0005 \\
    p99 $|\Delta|$ & 0.143 \\
    Max $|\Delta|$ & 0.422 \\
    Top-1 agreement & 0.991 \\
    Top-20 overlap & 0.985 \\
    \bottomrule
  \end{tabular}
\end{table}

% generated by scripts/make_paper_tables.py from runs/rollout_bench/*/summary.json; do not edit
\begin{table}[t]
  \centering\small
  \caption{Student rollout throughput (8 prompts $\times$ 2 samples, T=1, up to 4{,}096 tokens).}\label{tab:throughput}
  \begin{tabular}{lrrr}
    \toprule
    Student & tokens/s & mean length & engine load (s) \\
    \midrule
    S2 & 844 & 1{,}115 & 116 \\
    S1 & 376 & 1{,}438 & 124 \\
    \bottomrule
  \end{tabular}
\end{table}


\section{Quantization recipes}\label{app:quant}
All rungs are built with llm-compressor (0.15.0) from the same \calibN{} calibration texts, with
every mixture-of-experts expert calibrated and the router gates and output head kept in BF16:
FP8 W8A8 with static activation scales (T8), AWQ W4A16 with group size 128 (T4-int), NVFP4
(T4-fp) and GPTQ W3A16 with group size 128 (T3). Quantization runs in a separate container from
inference, and every checkpoint is loaded by the same vLLM build used for scoring. Before building
teacher rungs we checked each format on the 0.6B student (Table~\ref{tab:quantval}); the formats
order as expected, from FP8 closest to BF16 to 3-bit furthest.
\todo{per-rung tool versions, timings and calibration hash from each rung's \texttt{recipe.json}}

% generated by scripts/make_paper_tables.py from reports/quant_validation.json; do not edit
\begin{table}[t]
  \centering\small
  \caption{Quantization tooling check on the 0.6B student: each format's checkpoint, served by vLLM, against BF16 on the same 20 rollouts (log-probabilities of the realized tokens, top-20).}\label{tab:quantval}
  \begin{tabular}{lrrrrr}
    \toprule
    Format & mean $|\Delta|$ & median & p99 & top-1 agr. & top-20 overlap \\
    \midrule
    FP8 W8A8 (static) & 0.076 & 0.0063 & 0.77 & 0.961 & 0.932 \\
    AWQ W4A16 g128 & 0.272 & 0.0275 & 2.73 & 0.879 & 0.802 \\
    NVFP4 & 0.313 & 0.0397 & 3.06 & 0.857 & 0.782 \\
    GPTQ W3A16 g128 & 0.685 & 0.1082 & 6.20 & 0.767 & 0.637 \\
    \bottomrule
  \end{tabular}
\end{table}

 \todo{tool versions, calibration hash and settings per rung, from each rung's
\texttt{recipe.json}.}

\section{Per-benchmark results}\label{app:perbench}
\TBD{}

\section{Compute log}\label{app:compute}
\TBD{} \todo{machine-hours per experiment, from the run directories; compare with the
pre-registered budget.}
